
Practitioners who fine-tune code language models must decide which datasets to train on and in what proportion. The prevailing assumption is that training on problems similar in surface form to the evaluation benchmark produces the best results --- that MBPP-style utility scripts should improve MBPP performance, and that HumanEval-style algorithm problems should improve HumanEval performance. If this assumption were correct, data source selection would reduce to benchmark-matching.

Empirical evidence for or against this assumption is absent from the literature. Existing code SFT papers --- WizardCoder \citep{Luo2023WizardCoder}, OctoPack \citep{Muennighoff2023OctoPack}, DeepSeek-Coder \citep{Guo2024DeepSeekCoder} --- report training data compositions without ablating source identity. They vary data quality, quantity, and format simultaneously, making causal attribution of any performance difference to the source identity alone impossible. Domain mixture optimization literature [Xie et al., 2023; Xie et al., 2025; Zhang, 2026] addresses pretraining data mix but has not been applied as a controlled source-identity ablation for code-specific SFT with execution-based evaluation.

This paper closes that gap. We train DeepSeek-Coder at 1.3B and 7B scales on four fully isolated source conditions, holding token budget, hyperparameters, supervision format, and random seed constant across conditions. The result is a controlled causal estimate of how much source identity alone affects post-SFT pass@1.

The main findings are as follows. First, source identity produces large, statistically significant differences in HumanEval+ pass@1 at 1.3B scale: HumanEval-only achieves a mean of 35.0\% while LeetCode-only achieves ~3.1\% under identical training budgets (ANOVA F=11.37, p=0.020, max contrast 31.9pp). Second, the assumed symmetric specialization pattern does not hold: HumanEval-only SFT achieves higher MBPP+ pass@1 (~51.9\%) than MBPP-only SFT (~50.5\%) consistently across all three seeds, refuting the ''train on X to test on X'' principle in this direction. Third, CodeBERT embedding cosine similarity between training source and test benchmark rank-orders the four conditions in perfect agreement with their HumanEval+ pass@1 rank ($\rho$=1.0, p=0.042), providing evidence that distributional alignment in code-embedding space is a pre-training predictor of SFT behavioral rank. Fourth, at 7B scale, the condition rank order on HumanEval+ is preserved, but the absolute spread narrows from 31.9pp to 5.5pp, and equal-mix (39.6\%) matches or exceeds HumanEval-only (38.6\%), suggesting that larger-scale pretraining reduces the marginal impact of SFT source specificity.

This work contributes to the DL4C topics of data for code and post-training methods for code LLMs. The experimental controls --- token-budget equalization, cosine-similarity deduplication against test benchmarks, format normalization --- are designed to isolate source identity as the single independent variable.

